\section{相対Poisson臨界スキーム、二重五次判別式、\texorpdfstring{$S_5$}{S5}モノドロミー}
\label{sec:relative-poisson-cover}
前節の十Poisson零点を、低次係数$(a,b,c)$を動かした相対スキームとして調べる。
以下では被約supportとスキームの長さを区別する。基礎体$k$は標数0の代数閉体、
$q^2+q+1=0$、$f=(1-q)/3$、$e=-(2q+1)/9$とする。

\subsection{全係数で有限平坦な二つの被覆}
$A=k[a,b,c]$、$S=\Spec A$とし、式\eqref{eq:center}の左辺を
$P_{\mathbf a}(X,Y,Z,\Omega)$と書く。Casimir切断の特異点を添字付ける
相対Poisson臨界スキームは
\[
 \mathcal K=\Spec\frac{A[X,Y,Z,\Omega]}
 {(P_{\mathbf a},P_X,P_Y,P_Z)}\xrightarrow{\kappa}S
\]
である。ここで偏微分は$X,Y,Z$に関するもので、$\Omega$は固定しない。

\begin{theorem}[次数十五と五の全係数有限平坦性]
\label{thm:relative-critical-flat}
$\kappa$は有限平坦次数$15$である。また一次元表現の相対スキーム
\[
 \mathcal E=\Spec\frac{A[x,y,z]}
 {((1-q)yz-x-a,\ (1-q)zx-y-b,\ (1-q)xy-z-c)}
 \xrightarrow{\pi}S
\]
は有限平坦次数$5$であり、総空間は$\AA^3_{x,y,z}$と同型なので既約である。
$\pi$は原点を含む非空Zariski開集合上で有限\'{e}taleとなる。
\end{theorem}
\begin{proof}
$P_{\mathbf a}$を通常の全次数で斉次化し、
$[X:Y:Z:\Omega:T]\in\PP^4_S$で三次式$P^h$と
三二次式$P_X^h,P_Y^h,P_Z^h$の共通零点を$\overline{\mathcal K}$とする。
$T=0$では
\[
 (P^h,P_X^h,P_Y^h,P_Z^h)
 =(\Omega^3-XYZ,-YZ,-ZX,-XY),
\]
従ってsupportは三座標頂点のみである。
$X=1$の頂点近傍では
\begin{align*}
P_X^h&=-YZ+T(2e-a\Omega)+T^2\ell_a,\\
P_Y^h&=-Z+T(2eY-b\Omega)+T^2\ell_b,\\
P_Z^h&=-Y+T(2eZ-c\Omega)+T^2\ell_c .
\end{align*}
これらの$(T,Z,Y)$に関するJacobianは頂点で可逆（行列式$\pm2e$）。
$T=Y=Z=0$が$\Omega$の任意値に対する厳密解なので形式陰関数定理から、
三式で割った完備局所環では$T=Y=Z=0$、残る関係は$\Omega^3=0$となる。
よって各無限遠頂点は相対長さ3の自明な厚みである。

任意の正次元射影部分多様体は無限遠と交わるが、上の三頂点は孤立している。
従って全幾何学的fiberは零次元で、$\overline{\mathcal K}$は
次数$(3,2,2,2)$の相対完全交叉である。
B\'{e}zoutから各fiberの長さは$3\cdot2^3=24$。
射影かつ準有限だから有限であり、完全交叉のCohen--Macaulay性と
regularな底$S$へのmiracle flatnessにより平坦である。
無限遠の三つの厚みは開閉な有限平坦部分（次数9）なので、
その補集合である$\mathcal K$は有限平坦次数$24-9=15$となる。

一次元表現について三二次関係を$[x:y:z:T]\in\PP^3_S$に斉次化する。
$T=0$では三座標頂点のみで、例えば$x=1$で
\[
 (1-q)yz-T-aT^2,\quad (1-q)z-yT-bT^2,\quad
 (1-q)y-zT-cT^2
\]
の$(T,z,y)$に関するJacobianは単元$\pm(1-q)^2$である。
三つの無限遠点は長さ1の自明な開閉section。
射影三二次完全交叉は有限平坦次数$8$なので、これらを除く
$\mathcal E$は次数$8-3=5$となる。
$a,b,c$は$x,y,z$の多項式だから総空間は$\AA^3$。
原点fiberの五点のJacobianは$-1,4,4,4,4$で非零
（原点と $(x,y,z)=(\eta_1,\eta_2,\eta_3)/(1-q)$、$\eta_i=\pm1$、$\eta_1\eta_2\eta_3=1$ への直接代入）なので\'{e}tale開集合は非空である。
\end{proof}

\subsection{二重五次因子とコンパニオン五次因子}
各一次元表現に$X,Y,Z,\Omega$を評価すると、
全係数で$P=P_X=P_Y=P_Z=P_\Omega=0$が成立する。
これは元関係の代入恒等式であり、
\texttt{computations/verify\_character\_node\_cover.py}でも検証済みである。

\begin{theorem}[二重節点とCasimir臨界値の五次因数分解]
\label{thm:critical-discriminant-quintics}
原点を含む非空Zariski開近傍$S^\circ$上で、$\mathcal K$は
相対節点に沿う次数10の二重厚みと、総中心の滑らかなPoisson臨界点を
添字付ける有限\'{e}tale五重被覆$\mathcal E'\to S^\circ$に開閉分解する。
前者は$\mathcal E|_{S^\circ}$の相対一次無限小厚みで、
\'{e}tale局所的には
\[
 \mathcal K|_{S^\circ}\simeq
 \Spec_{\mathcal E}(\OO_{\mathcal E}[\epsilon]/\epsilon^2)
 \sqcup\mathcal E'.
\]
$\mathcal E'$の幾何学的既約性・monodromyは、後述の定理\ref{thm:paired-quintic-S5}で完全決定する。

二つのmonicな特性多項式
\[
 N_5(D)=\det(D+\Omega\mid\pi_*\OO_{\mathcal E}),\qquad
 \Delta_{15}(D)=\det(D+\Omega\mid\kappa_*\OO_{\mathcal K})
\]
は$A[D]$上の次数5と15の多項式であり、
\begin{equation}\label{eq:relative-discriminant-quintics}
 \boxed{\Delta_{15}(D)=N_5(D)^2M_5(D)},\qquad
 M_5\in A[D]\text{ monic of degree }5
\end{equation}
が\emph{全係数}で成立する。
$\Delta_{15}=0$はCasimir中心切断が特異となる値の条件、
$N_5=0$は一次元表現に由来する節点の中心値である。
原点での特殊化は
\[
 N_5(D)=D(D-2f/3)^4,\qquad
 M_5(D)=(D-f)(D+f/3)^4.
\]
\end{theorem}
\begin{proof}
原点fiberの臨界点supportは定理\ref{thm:ten-poisson-points}の十点。
全点で$(X,Y,Z)$-Hessianが非退化なので、$P_X=P_Y=P_Z=0$から
$X,Y,Z$を$(\Omega,a,b,c)$の形式関数に解ける。
これを$P$に代入した一変数方程式を$h(\Omega)$とする。
総中心の五節点では$P_\Omega=0$、$P_{\Omega\Omega}=2f,-2f$は単元。
一次元表現による節点sectionは全係数で
$P=P_X=P_Y=P_Z=P_\Omega=0$を満たして存続する。
従って各節点上で$h=(\Omega-\Omega_{\rm node})^2\cdot\mathrm{unit}$となる。
他の五点では$P_\Omega=f^2$が単元なので、四式のJacobianが可逆で、
相対長さ1の有限\'{e}tale点へ延長する。
有限平坦性と原点の十点の相異性から、底を開に縮めればこれらは
互いに素な開閉部分へ分離する。次数は$5\cdot2+5=15$である。

この分解における$D+\Omega$の特性多項式は
相対節点側で$N_5(D)^2$、残部でmonic五次式になる。
従って$S^\circ$上で式\eqref{eq:relative-discriminant-quintics}が成立する。
$A=k[a,b,c]$は整域で$N_5$はmonicなので、
$A[D]$で$\Delta_{15}$を$N_5^2$で割った剰余が
稠密開上で零なら全域で零となる。ゆえに$M_5\in A[D]$が一意に存在する。
原点での式は、五つの$(X,Y,Z)$-臨界三つ組に対する$\Omega$多項式が
一つは$\Omega^2(\Omega+f)$、残る四つは
$(\Omega-f/3)(\Omega+2f/3)^2$であることから従う。
\texttt{verify\_relative\_poisson\_cover.py}は
$\mathbb Q(\sqrt{-3})$上のGr\"obner商の長さ15、
$-\Omega$の$15\times15$乗法行列の特性多項式、
無限遠のJacobianを独立に厳密検算する。
\end{proof}

\subsection{既約五次因子の\texorpdfstring{$S_5$}{S5}モノドロミー}
\begin{theorem}[一次元表現被覆の完全対称群]
\label{thm:character-cover-S5}
$\mathcal E\to S$の最大\'{e}tale開部分の幾何学的monodromy群は
五文字の完全対称群$S_5$である。
さらに$N_5(D)$は$k(a,b,c)[D]$上既約・分離的で、
その分解体のGalois群も$S_5$である。
\end{theorem}
\begin{proof}
$r=1-q$、$u=rx,v=ry,w=rz$、$A=ra,B=rb,C=rc$と置くと、
一次元表現の係数写像は
\begin{equation}\label{eq:five-character-gradient}
 (u,v,w)\mapsto(A,B,C)=(vw-u,wu-v,uv-w)
\end{equation}
となる。これは$uvw-(u^2+v^2+w^2)/2$のgradient mapである。
被覆の総空間は既約な$\AA^3_{u,v,w}$なのでmonodromyは五点に推移的であり、
次数が素数5なのでprimitiveである。
Jacobian行列と行列式は
\[
 J=\begin{pmatrix}-1&w&v\\w&-1&u\\v&u&-1\end{pmatrix},
 \quad\det J=u^2+v^2+w^2+2uvw-1.
\]
$A=B=C=-1/4$のfiberには$(1/2,1/2,1/2)$が重複度2で現れ、
他は$(5/4,-1,-1)$の三順列（Jacobian行列式$81/16$）である。
前者で$J$の核・余核は対角方向で、対角写像
$t\mapsto t^2-t$は$t=1/2$で単純foldとなる。
従って幾何学的慣性群には一つの互換がある。
primitive置換群において一つの互換の全共役を辺とするグラフは
連結でなければ非自明なblock系を作る。
連結グラフの辺に沿う互換は$S_5$を生成するから、
monodromy群は$S_5$である。

一次元表現条件を式\eqref{eq:Omega}へ代入し
$q^2+q+1=0$で還元すると、
\begin{equation}\label{eq:five-character-casimir}
 \Omega_{\mathcal E}(u,v,w)=\frac f3(uvw-u^2-v^2-w^2)
\end{equation}
を得る。原点fiberの五点は$(0,0,0)$および
$\eta=(\eta_1,\eta_2,\eta_3)$、$\eta_i=\pm1$、積1である。
前者の$\Omega$値は$0$、後四点は$-2f/3$だが、
その底$(A,B,C)$上の一次微分は
\[
 d\Omega_\eta=-\frac f3(\eta_1dA+\eta_2dB+\eta_3dC).
\]
これは$\nabla_{u,v,w}\Omega_\eta=-f\eta/3$かつ$J(\eta)\eta=\eta$から従う。
四つの一次形式は相異なるので、五つの$\Omega$値は
一般の底上で相異なる。従って$\Omega_{\mathcal E}$は
次数5の拡大$k(\mathcal E)/k(a,b,c)$の原始元で、
$N_5$は既約・分離的、そのGalois群は先ほどの$S_5$である。
\texttt{verify\_relative\_poisson\_cover.py}はsimple fold、
Casimirの式と一次微分を全係数で厳密検算する。
\end{proof}

\subsection{二つの五次因子の独立な\texorpdfstring{$S_5$}{S5}モノドロミー}
\begin{theorem}[コンパニオン五次式の幾何学的$S_5$と独立性]
\label{thm:paired-quintic-S5}
式\eqref{eq:relative-discriminant-quintics}の二つの五次因子
$N_5(D),M_5(D)$について、標数零の任意の代数閉体$k$上で次が成り立つ。
\begin{enumerate}
\item $M_5(D)$は$k(a,b,c)[D]$上既約かつ分離的であり、
その幾何学的Galois群は$S_5$である。
従って定理\ref{thm:critical-discriminant-quintics}の
残余有限\'{e}tale五重被覆$\mathcal E'$は幾何学的に連結である。
\item $N_5,M_5$の分解体をそれぞれ$L_N,L_M$とすると
\[
 \operatorname{Gal}(L_NL_M/k(a,b,c))\cong S_5\times S_5.
\]
特に両分解体は線型非交叉で、二つの判別式が定める
$k(a,b,c)^\times/k(a,b,c)^{\times2}$の平方類も独立である。
\end{enumerate}
この結果は二つの被覆のGalois幾何を決めるが、
一般$(a,b,c)$の非可換H3性・全Azumaya locus・局所Morita型を
それだけで判定するものではない。
\end{theorem}
\begin{proof}
まず$K_0=\mathbb Q(q)$で考える。
全係数の有限平坦性は$\mathbb Z[q,1/6]$上でも成立し、
$q^2+q+1=0$を用いた係数表示の分母は$3$の冪なので、
以下の素数における有限体還元はよく定義される。
$K_0$-有理な係数$(a,b,c)=(30,19,29)$で特殊化する。
次の二つの根の選択は、それぞれ$K_0$の素イデアル
$q\equiv374\pmod{1009}$、$q\equiv86\pmod{1069}$を指定する。
両方で$q^2+q+1\equiv0$が成立する。

各有限体上で$P,P_X,P_Y,P_Z$のGr\"obner商の長さが15であることを
確認し、$-\Omega$乗法の$15\times15$行列を構成する。
これとは独立に、一次元表現の三二次関係のGr\"obner商の長さ5を確認し、
式\eqref{eq:five-character-casimir}の$-\Omega$乗法行列から
$N_5(D)$を得る。前者の特性多項式を$N_5(D)^2$で割ると
余り零、商は$M_5(D)$となる。正確な結果は次の通りである。
\begin{align*}
 M_5\bmod(1009,q-374)
  ={}&D^5-404D^4-215D^3+279D^2+226D+42,\\
 M_5\bmod(1069,q-86)
  ={}&(D^2+248D+337)(D^3+55D^2-276D-260).
\end{align*}
前者は$\mathbb F_{1009}$上既約、後者の二因子は
$\mathbb F_{1069}$上それぞれ既約で互いに異なる。
後者の素数では$N_5$は
\[
 D^5-100D^4-59D^3-17D^2+487D+534
\]
であり、こちらも$\mathbb F_{1069}$上既約で、$M_5$と互いに素である。
両素数で多項式の判別式は零でない。
これらの有限体上の等式、因数分解、既約性は
\texttt{computations/verify\_companion\_s5.py}で
元の中心と一次元表現の関係式から正確に再生成できる。

有限体のFrobeniusの分解型はそれぞれ$(5)$と$(2,3)$なので、
$K_0$上の特殊化$M_5$の算術Galois群には
位数5の元と型$(2,3)$の元がある。
後者は奇置換かつ位数6である。
従って群位数は$30$の倍数で、$S_5$の部分群だから
$30,60,120$のいずれかである。
位数60なら指数2の部分群$A_5$となり奇置換を含めない。
位数30なら$S_5$の指数4の部分群を与えるが、
四剰余類への推移的作用$S_5\to S_4$は存在しない
（$S_5$の固有非自明正規部分群は$A_5$のみ）。
よってGalois群は$S_5$である。
これを一般係数へ特殊化のGalois群の包含により戻すと、
$K_0(a,b,c)$上の$M_5$の算術Galois群も$S_5$である。

ここで算術的結論と幾何学的結論を区別する。
定理\ref{thm:critical-discriminant-quintics}の
$K_0$上の非空Zariski開集合$S^\circ$は係数原点を含み、
$\mathcal E'\to S^\circ$は有限\'{e}tale次数5である。
原点の五点は定理\ref{thm:ten-poisson-points}の
$r_0,r_\eta$で、すべて$K_0$-有理かつ相異なる。
従って$\mathcal E'$の原点fiberは$K_0^5$に完全分裂する。
$S^\circ$は幾何学的に連結であるから、\cite[Lemma 58.14.3]{StacksPi1}
の算術・幾何基本群完全列
\[
 1\longrightarrow\pi_1(S^\circ_{\overline K_0})
 \longrightarrow\pi_1(S^\circ)
 \longrightarrow\operatorname{Gal}(\overline K_0/K_0)\longrightarrow1
\]
を使う。原点という$K_0$-有理基点は右側へのsectionを与え、
$\mathcal E'$の五点fiberが完全分裂しているので
そのsectionの$S_5$への作用は自明である。
従って算術的な$S_5$作用は幾何基本群の像そのものであり、
$M_5$の幾何学的Galois群も$S_5$となる。
ここで$M_5(D)$自体は原点で四重根を持つが、
\emph{エタールな被覆は多項式の零点スキームではなく
$(X,Y,Z,\Omega)$を保持した$\mathcal E'$である}点に注意する。
これが(1)を示す。

(2)を示す。$N_5$側の算術・幾何群は
定理\ref{thm:character-cover-S5}により$S_5$であり、
$M_5$側も上で確定した。合併した算術monodromy
$G\subset S_5\times S_5$は両射影へ全射である。
$p=1069$では$N_5$の分解型が$(5)$で偶置換、
$M_5$の分解型が$(2,3)$で奇置換なので、
$G$には符号の異なる対が存在する。
Goursatの補題と$S_5$の正規部分群
$\{1\},A_5,S_5$から、真のsubdirect productは
二つの符号が一致する$C_2$上のfiber product、または
$S_5$の自己同型のグラフに限られる。
いずれも符号の異なる対を含まないから
$G=S_5\times S_5$である。

幾何学的同時monodromyについては、同じ有理基点における
$\mathcal E$の五点も$p_0,p_\eta$として全て$K_0$-有理である。
よって$\mathcal E\sqcup\mathcal E'$の十点fiberは完全分裂し、
基本群完全列の有理sectionは$S_5\times S_5$への作用が自明。
従って算術と幾何の同時monodromyは一致する。
特に合成分解体は$S_5\times S_5$をGalois群に持ち、
二分解体は線型非交叉である。
その二つの符号商$S_5\to C_2$は互いに独立だから、
判別式の平方類も独立となる。
$\overline K_0$から任意の標数零代数閉体$k$への拡大でも、
有限型Galois閉包の幾何学的連結性は保たれる。
\end{proof}

\begin{remark}[計算の意味と未解決事項]
二つの$S_5$の独立性は有限体の無作為標本からの外挿ではない。
全係数で定義したmonic五次式について、二つの指定した良い素数における
\emph{厳密な}因数分解からFrobeniusの置換型を証明し、
群論とエタール基本群の完全列を適用したものである。
これによって残余五次式$M_5$の幾何学的既約性・monodromyは解決した。
未決定なのは両五次式の具体的な大域係数の簡潔な表示、
分岐divisorの幾何学、局所order/H3判定との一般係数での関係、
および楕円$E_6$成分への移行である。
\end{remark}

\begin{remark}[Painlev\'{e} VIの判別式との区別]
Inaba--Iwasaki--Saito \cite[\S9.1, Definition 9.1]{IISDiscriminant}は
Fricke cubicの判別式を境界トレースの係数空間に引き戻すと
二乗因子が現れることを示す。
本稿の$N_5^2M_5$は\emph{別のパラメータ写像}の
次数十五の臨界値多項式であり、両者の因子が等しいとは主張しない。
五次式$M_5$の既約性・幾何学的$S_5$および$N_5$との独立性は
定理\ref{thm:paired-quintic-S5}で解決した。
一般$(a,b,c)$でのH3条件、局所order、楕円境界の八次元成分への
相対中心の延長と分岐divisorの比較は未解決である。
本稿の計算は既知のPainlev\'{e}判別式の再発見とは区別し、
個別の新規性は一次文献との詳細照合を待つ。
\end{remark}
